% AI-produced by GPT-6 Astra (OpenAI, ChatGPT Pro mode). No priority claim.
\subsection{A full ordered predimension category}

For a finite simple $3$-uniform hypergraph $X$ and a vertex subset $U$,
put
\[
 \delta_X(U)=2|U|-3|E_X(U)|,
\]
where $E_X(U)$ consists of the hyperedges contained in $U$.
Let $\mathcal C_\delta$ be a small skeleton of the category of nonempty
finite linearly ordered such hypergraphs satisfying
$\delta_X(U)>0$ for every nonempty $U$. Its arrows are increasing
induced embeddings with \emph{$d$-closed} image, where
\[
 A\leq_d B
 \quad\Longleftrightarrow\quad
 \delta_B(U)>\delta_B(A)
 \quad\text{for every }A\subsetneq U\subseteq B.
 \tag{DC}\label{eq:predimension-dclosed}
\]
Isomorphisms are allowed. Including the empty object instead makes
no difference to the assertions below. The strict closure terminology
is used for sparse graphs in
\citet[Definition~4.4]{EvansHubickaNesetril2019}; the weighted
hypergraph assertions here will be proved directly.

\begin{lemma}[Proved: closure calculus and free amalgamation]
\label{A:predimension-amalgam}
The relation $\leq_d$ is preserved by restriction to subsets, is
transitive, and is closed under finite intersections inside an ambient
object. Free amalgams of $d$-closed embeddings are positive-predimension
hypergraphs, and both factors are $d$-closed in the amalgam. The orders
can be amalgamated. In particular, $\mathcal C_\delta$ satisfies
condition I and has common successors.
\end{lemma}
\begin{proof}
Predimension is submodular:
\[
 \delta(U)+\delta(V)\geq\delta(U\cup V)+\delta(U\cap V).
 \tag{SM}\label{eq:predimension-submodular}
\]
The vertex term is modular, and the edge count is supermodular.
If $A\leq_d B$, $X\subseteq B$, and
$A\cap X\subsetneq U\subseteq X$, then $U\cap A=A\cap X$ and
$U\cup A\supsetneq A$. Hence
\[
 \delta(U)-\delta(A\cap X)
 \geq\delta(U\cup A)-\delta(A)>0.
\]
Thus $A\cap X\leq_d X$. For $A\leq_d B\leq_d C$ and
$A\subsetneq U\subseteq C$, the case $U\subseteq B$ is immediate.
Otherwise restriction gives $U\cap B\leq_d U$ and
\[
 \delta(U)>\delta(U\cap B)\geq\delta(A).
\]
This proves transitivity. Intersections follow by restriction and
transitivity.

Let $A\leq_d B,C$ and let $D=B\cup_A C$ be the free hypergraph
amalgam. For every $U\subseteq D$,
\[
 \delta_D(U)=\delta_B(U\cap B)+\delta_C(U\cap C)
              -\delta_A(U\cap A).
\]
If $U$ meets $C\setminus A$, restriction makes the difference of
the last two terms positive. The first is nonnegative. If $U$ does
not meet $C\setminus A$, it is a subset of $B$. Therefore every
nonempty $U$ has positive predimension. If $B\subsetneq U\subseteq D$,
then
\[
 \delta_D(U)-\delta(B)=\delta_C(U\cap C)-\delta(A)>0.
\]
Thus $B\leq_d D$, and symmetrically $C\leq_d D$. Merge the two
linear orders separately in each interval determined by their common
ordered subset $A$. This extends both orders. Amalgamation gives I;
the same construction over the empty structure gives common
successors, even if the empty structure is not an object.
\end{proof}

\begin{lemma}[Proved: the long-necklace predimension estimate]
\label{A:predimension-necklace}
Let $f,g:S\rightrightarrows T$ be arrows of $\mathcal C_\delta$.
Choose $n\geq2$ bounding the vertex lengths of the nonfixed paths
of the partial bijection $f(x)\mapsto g(x)$, and put $w=2n-2$.
For $k\geq n$, let $Z$ be the covered, actual-length repaired point
necklace of Lemma~\ref{C:necklace-set-colimit}, of length $s=4k+1$,
with copies $\alpha_i:T\to Z$. Give $Z$ precisely the hyperedges
that occur in those copies. If
\[
 s>w\bigl(\delta(T)+2|S|\bigr),
 \tag{PD}\label{eq:predimension-necklace-bound}
\]
then $Z$ is a positive-predimension hypergraph and every
$\alpha_i$ is an induced embedding with $d$-closed image.
\end{lemma}
\begin{proof}
We use the exact consecutive overlaps and visibility intervals from
Lemma~\ref{C:necklace-set-colimit}. A fixed point has full-circle
visibility. Every other point, including a private point, has a proper
cyclic interval of at most $w$ indices. Since $w\leq2k-2$, every
proper interval is shorter than half the index circle and contains
the unique shortest index path between any two of its indices.

First, each copy is induced. If an edge from copy $j$ has all its
vertices in copy $i$, the shortest index path from $j$ to $i$ lies
in all three vertex visibility sets. At a consecutive overlap the
edge is contained in the indicated $S$-image. Because $f$ and $g$
are induced embeddings, it comes from an edge of $S$ and propagates
to the next copy. It therefore occurs in copy $i$. This also shows
that the visibility of an edge is exactly the intersection of the
three vertex visibility sets. That intersection is a connected proper
interval, or the full circle when all three vertices are fixed.
All edge and point identifications are consequently generated by
the neighboring overlaps.

Cut one overlap of the necklace. Its linear amalgam $L$ is obtained
by successive free amalgamations over $S$. By
Lemma~\ref{A:predimension-amalgam}, it is positive-predimension,
and every $T$-copy is $d$-closed in it. Restoring the omitted overlap
gives a quotient $q:L\to Z$ which identifies at most $|S|$ vertex
pairs. If $U\subseteq Z$ and $\widetilde U=q^{-1}(U)$, every edge
of $Z[U]$ lifts to an edge of $L[\widetilde U]$. Hence
\[
 \delta_Z(U)\geq\delta_L(\widetilde U)-2|S|.
 \tag{1}\label{eq:predimension-cut-loss}
\]
Indeed, if $d_v,d_e$ are the numbers of vertex and edge duplications
removed by $q$ on these induced substructures, then the exact
difference is $2d_v-3d_e$, with $0\leq d_v\leq|S|$ and $d_e\geq0$.

Suppose first that the union of the proper visibility intervals of
the vertices in $U$ misses an index $r$. Cut the overlap from $r-1$
to $r$. No proper vertex fiber of $U$ is cut into two components,
and a full-circle fiber becomes a connected path. All such vertices
therefore have unique lifts. Edges have unique lifts as well, either
by the connected visibility description or by simplicity of the
hypergraph once their three vertices lift uniquely. Thus
$L[\widetilde U]\cong Z[U]$. This proves $\delta_Z(U)>0$ for
nonempty $U$. If $U$ properly contains a copy $\alpha_i(T)$, its
lift properly contains that same $T$-copy, so
$\delta_Z(U)>\delta(T)$.

It remains to consider the case where those proper visibility
intervals cover all $s$ indices. Cut any overlap and number the
linear copies $0,\ldots,s-1$. Build $L[\widetilde U]$ by successively
adjoining its intersection $U_j$ with copy $T_j$. For $j>0$, the
intersection of $T_j$ with the preceding amalgam is exactly its
incoming $S$-image $A_j$. Its predimension contribution is therefore
\[
 \delta(U_j)-\delta(U_j\cap A_j).
\]
This is zero if $U_j\subseteq A_j$. Otherwise submodularity and
$A_j\leq_d T_j$ give
\[
 \delta(U_j)-\delta(U_j\cap A_j)
 \geq\delta(U_j\cup A_j)-\delta(A_j)\geq1.
 \tag{2}\label{eq:predimension-new-stage}
\]
The initial contribution $\delta(U_0)$ is at least one whenever
$U_0$ is nonempty.

Let $J$ be the set of indices at which at least one nonfixed vertex
of $\widetilde U$ first appears. Each such vertex has a linear
visibility interval of at most $w$ indices. Consequently all the
nonfixed vertices first appearing at $j\in J$ are visible only
among $j,\ldots,j+w-1$. Every copy index contains a nonfixed vertex
of $U$ by the case assumption, and its lift has some first-appearance
index in $J$. The intervals $[j,j+w-1]$, for $j\in J$, therefore
cover all $s$ indices. Thus $|J|\geq\lceil s/w\rceil$. Every such
stage contributes at least one in
\eqref{eq:predimension-new-stage}, or in the initial contribution,
and all other stages contribute nonnegatively. We obtain
\[
 \delta_L(\widetilde U)\geq\lceil s/w\rceil,
 \qquad
 \delta_Z(U)\geq\lceil s/w\rceil-2|S|>\delta(T)>0,
\]
using \eqref{eq:predimension-cut-loss} and
\eqref{eq:predimension-necklace-bound}. This counts first-appearance
indices, not individual new vertices; no bound on $|T|$ is needed.
The two cases prove positivity for every nonempty subset and strict
relative positivity over every copy of $T$.
\end{proof}

\begin{proposition}[Proved: quasi-filteredness of the full category]
\label{A:predimension-quasi}
The category $\mathcal C_\delta$ is quasi-filtered. The same is true
of its full subcategory consisting of objects in which every point
and every hyperedge is $d$-closed.
\end{proposition}
\begin{proof}
For $f,g:S\rightrightarrows T$, choose $k$ satisfying
Lemma~\ref{A:predimension-necklace} and also exceeding the eventual
ordering threshold of Proposition~\ref{C:ordered-necklace}. The
latter proposition orders the very same point target $Z$ so that
every $\alpha_i$ is increasing. Hence every $\alpha_i$ is an arrow
in $\mathcal C_\delta$.

In the integral linearization, set
\[
 t=\sum_{i=0}^{2k}\alpha_i-\sum_{i=2k+1}^{4k}\alpha_i.
\]
It has augmentation one. The exact necklace equations give
\[
 tg=\sum_{i=1}^{2k}\alpha_i f
       -\sum_{i=2k+1}^{4k-1}\alpha_i f
    =tf-\alpha_0f+\alpha_{4k}f=tf.
\]
This is IV. Condition I and connectedness were proved in
Lemma~\ref{A:predimension-amalgam}.

For the asserted full subcategory, free amalgams and the necklace
targets are covered by old copies. Every new point or edge is
$d$-closed in some old copy, hence in the target by transitivity.
Both constructions therefore remain in this subcategory.
\end{proof}

\begin{proposition}[Proved: the initial triple remains unequalized]
\label{A:predimension-triple}
Let $V$ be one point, let $W$ be the ordered three-point single-edge
hypergraph, and let $v_1,v_2,v_3:V\to W$ be its three point arrows.
In either category of Proposition~\ref{A:predimension-quasi}, this
triple has no affine post-equalizer over any nonzero unital ring $R$.
For
\[
 D:R[A(W,-)]\longrightarrow R[A(V,-)]^2,
 \qquad D(g)=(gv_2-gv_1,gv_3-gv_1),
 \qquad Q=\operatorname{coker}D,
\]
one has $T_1^A(Q)\cong R$. In particular neither category is exact.
\end{proposition}
\begin{proof}
Each $v_i$ is $d$-closed: its predimension is two and its proper
extensions inside $W$ have predimension four or three. Consider a
nonzero coefficient vector supported on $r$ distinct arrows $g:W\to X$.
Their images are $r$ distinct hyperedges, since an increasing
three-point embedding is determined by its image. If their union has
$v$ vertices, positivity gives
\[
 2v>3|E_X(U)|\geq3r.
\]
The $r$ support edges have $3r$ incidences, so some vertex belongs
to exactly one support edge. Its coefficient $c\neq0$ occurs in
$D$ as $(-c,-c)$, $(c,0)$, or $(0,c)$ and cannot cancel. Its
corresponding $V$-arrow exists by composition $V\to W\to X$.
Thus $D$ is objectwise injective over every $R$.

The two displayed representable modules are projective. Their
colimits are $R$ and $R^2$, and the colimit of $D$ is zero because
each difference of two composites has augmentation zero. The long
exact derived-colimit sequence of this projective presentation gives
$T_1(Q)\cong R$. Objectwise injectivity of $D$ also directly
precludes any nonzero, hence any augmentation-one, simultaneous
equalizer of the three arrows.
\end{proof}

\paragraph{Scope and the subsequent totality proof.}
The preceding arguments establish quasi-filteredness and the
all-ring triple obstruction. The separate finite-support and
all-action arguments below prove totality of both categories in
Theorem~\ref{A:predimension-total}, giving the stated negative
answer to A.1.

In the unrestricted category, a raw vertex or hyperedge need not
be $d$-closed, so representable $V$- and $W$-actions cannot silently
be identified with all vertices and hyperedges. The preceding
argument uses only actual arrows. Also, the full category has
additional old-target arrows and additional amalgams. The selected
AKR class does not automatically survive upon passing to it.

